% ============================================================
% PART V: WRITTEN / FREE RESPONSE
%
% PEANUT EXAM WRITTEN QUESTION LAYOUT STANDARD
%
% AI INSTRUCTION:
% - Use \subsection*{Question 1}, \subsection*{Question 2}, ...
% - Sub-parts use enumerate with label=(\alph*)
% - Answer space is created with \vspace{...} ONLY. Never use \\[mm]
% - Size the space to the expected answer:
%     one-line recall        ~20mm
%     short explanation      ~30-40mm
%     derivation / multi-step ~45-60mm
%     full essay answer      ~80-90mm
% - Start a new \newpage before a question that will not fit on the page
% - State explicitly what a complete answer must contain
%   ("Your answer must include: ... , ... , and ...")
% - Do NOT alter math formatting; display math uses $$ ... $$
% ============================================================

\section*{Part V: Written Questions}

% -----------------------------------------------------------
\subsection*{Question 1}
% -----------------------------------------------------------

\begin{enumerate}[label=(\alph*)]

\item A conceptual question asking the student to explain a mechanism
in their own words. Your answer must include: the role of each component,
the order of the steps, and the property the mechanism guarantees.

\vspace{40mm}

\item A numeric or applied scenario. Given $x = 15$, $y = 4$, and $z = 12$,
state the conclusion and show the reasoning that leads to it.

\vspace{40mm}

\item A ``why'' question: explain what would break if one element were removed
(a nonce, a lock, a normalization step, a base case).

\vspace{45mm}

\item A short comparison: choose the right tool from a given list for two
different situations, and justify each choice in one or two sentences.

\vspace{30mm}

\end{enumerate}

% -----------------------------------------------------------
\newpage
\subsection*{Question 2}
% -----------------------------------------------------------

A question introducing given data, a matrix, a diagram, or a model equation.
For example, consider the approximation model
$$
y = c_1 x + c_2 x^2.
$$

\begin{enumerate}[label=(\alph*)]

\item Perform the computation or derivation. Show all work.

\vspace{60mm}

\item Interpret the result: what does it mean for the system in the scenario,
and what assumption would invalidate it?

\vspace{50mm}

\end{enumerate}

% -----------------------------------------------------------
\newpage
\subsection*{Question 3}
% -----------------------------------------------------------

A multi-step problem that integrates two topics from different lectures.
This is the question that separates memorization from understanding.

\begin{enumerate}[label=(\alph*)]

\item Step one of the analysis.

\vspace{50mm}

\item Step two, which depends on the answer to (a).

\vspace{50mm}

\item A closing judgment call: recommend an option and defend it.

\vspace{45mm}

\end{enumerate}
